%! Representing a Categorical Variable — Unit 1, Lesson 1.2 — Guided Notes & Practice (Answer Key)
% MIRROR of ../notes/main.tex: -key in place of -boxes, \termblank -> \termans, \blank -> \ans, and the
% answer argument of every \pcell, \writespace and \labelbox filled. Nothing else differs.
% THE in-class document, in the MAIN IDEAS / NOTES shape (LESSON_SHAPE.md §1): the vocabulary
% box, then ONE two-column guidednotes table. Each row is one idea — a label on the left; on the
% right a SHORT printed definition, ONE large pre-drawn display, then a grid of a few numbered
% problems with real writing room. Problems are numbered 1..N through the component.
% DENSITY: a \blank{} only where a word or number IS the answer (the share column of the
% frequency table); no mid-sentence blanks; two problems across; 1.8–2.6cm of answer space each.
%   I DO   : the teacher reads the definition, marks up the display, works the first problem.
%   WE DO  : the class works the rest of each grid, students holding the pen; the last row,
%            Guided Practice, is worked entirely together in a fresh context.
%   YOU DO : nothing on this page — ../homework/ IS the individual practice.
% The crux is the WINS FIRST row — South has more walkers AND North has the higher walking rate,
% both true, and the bigger group wins on every raw count before anyone counts anything
% (EK UNC-1.E.1). The plurality trap (problem 4) and the wrong-total trap (problem 15) are
% problems in a grid, not callouts; the crux is problem 12.
%
% This lesson absorbs TWO CED topics (1.3 tables and 1.4 graphs), so the Bar Graphs row keeps
% the mechanics to one pre-drawn chart and rides the axis idea into the crux row.
%
% THE DATA: North High, all 400 students — Bus 180, Car 120, Walk 80, Bike 20.
%           South High, all 1,000 students — Walk 150, Bike 20 (the district's side-by-side chart).
%           Two Shifts (guided practice) — Morning 48/108/84 of 240; Evening 18/30/32 of 80.
%
% PAGE LOCKSTEP with notes_key/: \termblank -> \termans, \blank -> \ans, and the answer argument
% of every \pcell, \writespace and \labelbox filled. Nothing else differs. A row cannot break
% across pages: figure and grid are separate sub-rows (`\\` then `& ...`).
\documentclass[12pt]{article}
\usepackage{apstats-article}
\usepackage{apstats-key}   % pulls in -boxes; do NOT also load -boxes

\begin{document}

\pageheader{Unit 1, Lesson 1.2}{Guided Notes \& Practice --- Answer Key}
% NAMESTRIP: no \namedateperiod here — the name row lives on the cover only.

\vspace{2pt}

\boxguard[16]
\begin{vocabbox}
\small
Fill in each term \textbf{as we name it} in the notes below.
\par
\termans{Frequency table}{a table of the number of cases in each category; the counts add to the group size.}
\par
\termans{Relative frequency table}{a table of the share of the group in each category; the entries always add to 1, or 100\%.}
\par
\termans{Bar graph}{one bar per category, its height the count or the relative frequency; gaps separate the bars.}
\par
\termans{The comparison rule}{to compare groups of different sizes, use relative frequencies --- never raw counts.}
\par
\end{vocabbox}

\small
\begin{guidednotes}
% ── ROW 1: frequency and relative frequency tables ───────────────────────────
\mainidea[Counts \&]{Shares} &
A \textbf{frequency table} lists the count in each category; the counts add to the group
size. A \textbf{relative frequency table} lists each category's share; the shares add to 1.

\vspace{5pt}
North High asked all \textbf{400} of its students how they get to school.

\vspace{3pt}
\textbf{1.} Fill the share column.

\vspace{3pt}
\begin{center}
\renewcommand{\arraystretch}{1.35}
\begin{tabular}{@{} l c c @{}}
 & a \labelbox{1.7cm}{count} & a \labelbox{1.7cm}{share} \\[3pt]
\rowcolor{sky}
\textbf{How they get to school} & \textbf{Frequency} & \textbf{Relative frequency} \\
\hline
Bus  & 180 & \ans{0.45 \quad (45\%)} \\
Car  & 120 & \ans{0.30 \quad (30\%)} \\
Walk &  80 & \ans{0.20 \quad (20\%)} \\
Bike &  20 & \ans{0.05 \quad (5\%)} \\
\hline
\textbf{Total} & \textbf{400} & \ans{1.00 \quad (100\%)} \\
\end{tabular}
\end{center}
\\
& \notesprompt{Read your table.}
\begin{probgrid}{|Y|Y|}
\pcell{2}{What do the four counts add to, and why must they?}{1.8cm}{400, the group size: each student is in exactly one category} &
\pcell{3}{Write 80 of 400 as a fraction, a decimal, a percent, and a rate.}{1.8cm}{$80/400 = 0.20 = 20\%$ = 1 in 5} \\ \hline
\end{probgrid}
\\ \hline
% ── ROW 2: plurality vs majority, and the context habit ──────────────────────
\mainidea[``Most'' is a promise]{About half} &
The largest category is a \textbf{plurality}. A \textbf{majority} is more than half. On the AP
exam a share always names its group: \emph{20\% of North students walk}.

\vspace{4pt}
\begin{center}
\begin{tikzpicture}[scale=1.0]
  \fill[navy]     (0,0)   rectangle (4.5,1.0);
  \fill[goldacc]  (4.5,0) rectangle (7.5,1.0);
  \fill[navy!35]  (7.5,0) rectangle (9.5,1.0);
  \fill[goldacc!45] (9.5,0) rectangle (10,1.0);
  \draw[white, thick] (4.5,0) -- (4.5,1.0);
  \draw[white, thick] (7.5,0) -- (7.5,1.0);
  \draw[white, thick] (9.5,0) -- (9.5,1.0);
  \draw (0,0) rectangle (10,1.0);
  \node[white, font=\small\bfseries] at (2.25,0.5) {Bus 45\%};
  \node[white, font=\small\bfseries] at (6.0,0.5) {Car 30\%};
  \node[font=\small\bfseries] at (8.5,0.5) {Walk 20\%};
  \node[font=\tiny, above] at (9.75,1.05) {Bike 5\%};
  \draw[redacc, very thick, dashed] (5,-0.45) -- (5,1.45);
  \node[redacc, font=\small\bfseries, below] at (5,-0.45) {half of North};
  \foreach \x/\l in {0/0\%, 2.5/25\%, 5/50\%, 7.5/75\%, 10/100\%} {
    \draw (\x,-0.05) -- (\x,-0.15); \node[below, font=\tiny] at (\x,-0.15) {\l};}
\end{tikzpicture}
\end{center}
\\
& \notesprompt{Check the claim.}
\begin{probgrid}{|Y|Y|}
\pcell{4}{A flyer says ``Most North students take the bus.'' Is that a majority?}{2.2cm}{45\% --- a plurality, not a majority: the other 55\% do something else} &
\pcell{5}{Rewrite ``20\%'' as a sentence an AP reader would accept.}{2.2cm}{20\% of North students walk to school} \\ \hline
\end{probgrid}
\\ \hline
% ── ROW 3: bar graphs ────────────────────────────────────────────────────────
\mainidea{Bar graphs} &
A \textbf{bar graph} draws the table: one bar per category, its height the count --- or the
share, if the axis is relabelled. Gaps separate the bars: categories are not a number line.

\vspace{4pt}
\begin{center}
\begin{tikzpicture}[scale=1.0]
  \foreach \y in {1,2,3} {\draw[linegray, very thin] (0,\y) -- (5.4,\y);}
  \draw[->, thick] (0,0) -- (0,4.1);
  \draw[->, thick] (0,0) -- (5.6,0);
  \foreach \y/\lab in {1/50, 2/100, 3/150} {
    \draw (-0.08,\y) -- (0.08,\y) node[left, font=\scriptsize, xshift=-2pt] {\lab};}
  \foreach \x/\h/\lab in {0.5/3.6/Bus, 1.8/2.4/Car, 3.1/1.6/Walk, 4.4/0.4/Bike} {
    \fill[navy] (\x,0) rectangle ++(0.8,\h);
    \node[below, font=\scriptsize] at (\x+0.4,-0.05) {\lab};}
  \node[rotate=90, font=\scriptsize] at (-0.95,2.0) {Students};
  \node[font=\scriptsize] at (2.8,-0.75) {How they get to school};
\end{tikzpicture}\\[6pt]
Each bar is one \labelbox{2.0cm}{category} \qquad The vertical axis is the \labelbox{2.0cm}{count}
\end{center}
\\
& \notesprompt{Read the graph.}
\begin{probgrid}{|Y|Y|}
\pcell{6}{How tall is the Walk bar, and what share of North is that?}{1.8cm}{80 students; $80/400 = 20\%$} &
\pcell{7}{Relabel the axis 0.45, 0.30, 0.20, 0.05. Which bars change height?}{1.8cm}{None --- only the axis labels change} \\ \hline
\end{probgrid}
\\
& \begin{probgrid}*{|Y|Y|}
\pcell{8}{Why are there gaps between the bars?}{1.8cm}{Categories are separate groups, not a number line} &
\pcell{9}{Before reading any bar graph, what is the first question to ask?}{1.8cm}{What is the vertical axis counting?} \\ \hline
\end{probgrid}
\\ \hline
% ── ROW 4: the comparison rule — THE CRUX ────────────────────────────────────
\mainidea[The bigger group]{Wins first} &
South High, in the same district, surveyed all \textbf{1{,}000} of its students --- 2.5~times
as many as North. The district posted this chart.

\vspace{3pt}
\begin{minipage}[c]{0.40\linewidth}
\renewcommand{\arraystretch}{1.3}\setlength{\tabcolsep}{4pt}
\begin{tabular}{@{} l r r r @{}}
\rowcolor{sky}
\textbf{School} & \textbf{Walk} & \textbf{Bike} & \textbf{Total} \\
\hline
North &  80 & 20 &   400 \\
South & 150 & 20 & 1{,}000 \\
\hline
\end{tabular}
\end{minipage}
\hfill
\begin{minipage}[c]{0.58\linewidth}
\centering
\begin{tikzpicture}[scale=0.85]
  \foreach \y in {1,2,3} {\draw[linegray, very thin] (0,\y) -- (6.6,\y);}
  \draw[->, thick] (0,0) -- (0,3.9);
  \draw[->, thick] (0,0) -- (6.8,0);
  \foreach \y/\lab in {1/200, 2/400, 3/600} {
    \draw (-0.08,\y) -- (0.08,\y) node[left, font=\scriptsize, xshift=-2pt] {\lab};}
  \foreach \x/\hn/\hs/\lab in {0.4/0.9/2.5/Bus, 2.0/0.6/1.65/Car, 3.6/0.4/0.75/Walk,
                               5.2/0.1/0.1/Bike} {
    \fill[navy]   (\x,0)     rectangle ++(0.55,\hn);
    \fill[goldacc](\x+0.6,0) rectangle ++(0.55,\hs);
    \node[below, font=\scriptsize] at (\x+0.57,-0.05) {\lab};}
  \fill[navy] (5.4,3.4) rectangle ++(0.3,0.3);
  \node[right, font=\scriptsize] at (5.72,3.55) {North};
  \fill[goldacc] (5.4,2.95) rectangle ++(0.3,0.3);
  \node[right, font=\scriptsize] at (5.72,3.1) {South};
  \node[rotate=90, font=\scriptsize] at (-1.0,1.9) {Students};
\end{tikzpicture}
\end{minipage}
\\
& \textbf{The comparison rule:} compare groups of different sizes with relative frequencies,
never raw counts. But ``how many?'' wants a count.

\vspace{3pt}
\textbf{In your own words:} why was South's bar guaranteed to be taller?
\writespace{1.6cm}{South asked 2.5 times as many students, so it wins on every raw count.}
\\
& \notesprompt{Compare the schools.}
\begin{probgrid}{|Y|Y|}
\pcell{10}{Work both walking rates.}{2.0cm}{North $80/400 = 20\%$; South $150/1000 = 15\%$} &
\pcell{11}{Which school has more walkers? Which has the higher walking rate?}{2.0cm}{South, 150 to 80; North, 20\% to 15\%} \\ \hline
\end{probgrid}
\\
& \begin{probgrid}*{|Y|Y|}
\pcell{12}{So which school ``walks more''? Answer honestly.}{2.2cm}{It depends on the question --- both statements are true} &
\pcell{13}{What one change would make the district's chart a fair comparison?}{2.2cm}{Relabel the axis as the percent of each school} \\ \hline
\end{probgrid}
\\
& \begin{probgrid}*{|Y|Y|}
\pcell{14}{``How many students in this district walk to school?'' Which number, and why?}{2.2cm}{A count, $80 + 150 = 230$: the question asks how many} &
\pcell{15}{A student divides South's 150 walkers by 1{,}400. What did she compute?}{2.2cm}{150 as a share of both schools --- the wrong total} \\ \hline
\end{probgrid}
\\ \hline
% ── ROW 5: guided practice — WE DO, together, fresh context ──────────────────
\mainidea[Guided practice]{Two Shifts} &
A coffee shop records every drink's size. Morning shift: \textbf{240} orders. Evening:
\textbf{80}.

\vspace{3pt}
\begin{minipage}[c]{0.44\linewidth}
\renewcommand{\arraystretch}{1.3}\setlength{\tabcolsep}{4pt}
\begin{tabular}{@{} l r r r r @{}}
\rowcolor{sky}
\textbf{Shift} & \textbf{S} & \textbf{M} & \textbf{L} & \textbf{Total} \\
\hline
Morning &  48 & 108 & 84 & 240 \\
Evening &  18 &  30 & 32 &  80 \\
\hline
\end{tabular}
\end{minipage}
\hfill
\begin{minipage}[c]{0.54\linewidth}
\centering
\begin{tikzpicture}[scale=0.8]
  \foreach \y in {1,2,3} {\draw[linegray, very thin] (0,\y) -- (5.0,\y);}
  \draw[->, thick] (0,0) -- (0,3.7);
  \draw[->, thick] (0,0) -- (5.2,0);
  \foreach \y/\lab in {1/40, 2/80, 3/120} {
    \draw (-0.08,\y) -- (0.08,\y) node[left, font=\scriptsize, xshift=-2pt] {\lab};}
  \foreach \x/\hm/\he/\lab in {0.4/1.2/0.45/Small, 2.0/2.7/0.75/Medium, 3.6/2.1/0.8/Large} {
    \fill[navy]   (\x,0)     rectangle ++(0.55,\hm);
    \fill[goldacc](\x+0.6,0) rectangle ++(0.55,\he);
    \node[below, font=\scriptsize] at (\x+0.57,-0.05) {\lab};}
  \fill[navy] (3.9,3.3) rectangle ++(0.3,0.3);
  \node[right, font=\scriptsize] at (4.22,3.45) {Morning};
  \fill[goldacc] (3.9,2.85) rectangle ++(0.3,0.3);
  \node[right, font=\scriptsize] at (4.22,3.0) {Evening};
  \node[rotate=90, font=\scriptsize] at (-1.0,1.8) {Orders};
\end{tikzpicture}
\end{minipage}
\\
& \notesprompt{We work these together.}
\begin{probgrid}{|Y|Y|}
\pcell{16}{Which shift sold more large drinks?}{2.0cm}{Morning, 84 to 32} &
\pcell{17}{Work both large-drink rates. Which shift sells them more often?}{2.0cm}{Morning $84/240 = 35\%$; evening $32/80 = 40\%$. Evening} \\ \hline
\end{probgrid}
\\
& \begin{probgrid}*{|Y|Y|}
\pcell{18}{One sentence for the staff newsletter comparing the shifts. Which kind of number
did you use, and why?}{2.6cm}{A rate: evening 40\% vs.\ 35\%, because 80 vs.\ 240 orders is not a fair count} &
\pcell{19}{The owner asks, ``Which shift is \emph{busier}?'' Which column answers that --- a
rate or a count?}{2.6cm}{The Total column: morning, 240 to 80 --- a count} \\ \hline
\end{probgrid}
\\ \hline
\end{guidednotes}

% ── THE NOTES END AT GUIDED PRACTICE ─────────────────────────────────────────
% There is deliberately no "you do" here: ../homework/ is the individual practice,
% started in class in the last 8 minutes while the teacher circulates.

\end{document}
